% Generated by research/polis/paper_examples.py -- do not edit by hand.
% Every Japanese string below is verbatim from the shipped corpus; the English is an
% author translation for readability only.
\appendix
\section{Training and Evaluation Instances}
\label{app:examples}

All instances below are verbatim from the corpus the experiments consume. English
renderings in italics are our translations, supplied for readability; no experiment
sees them. Speech identifiers are given so each instance can be located in the Diet
record.

\newcommand{\dexlabel}[1]{\par\vspace{3pt}\noindent{\footnotesize\scshape #1}\par\nobreak\vspace{1pt}}
\newcommand{\dexgloss}[1]{\par\vspace{1pt}\noindent{\footnotesize\itshape #1}\par}
\newenvironment{dexample}{%
  \par\vspace{4pt}\begin{list}{}{%
    \setlength{\leftmargin}{1.1em}\setlength{\rightmargin}{0em}%
    \setlength{\itemindent}{0em}\setlength{\parsep}{2pt}%
    \setlength{\topsep}{2pt}}\item[]\small%
}{\end{list}\vspace{2pt}}
\subsection{An anchor training instance (Section~\ref{sec:main}, ``Anchor Alignment'')}

Anchor \textbf{塩川鉄也} (日本共産党),
sitting of 2019-11-06, speech \texttt{120004889X00420191106\_121}. The
\emph{prompt} is the reconstructed prior context of the same sitting plus a speaker
cue; the \emph{chosen} completion is what the politician actually said next; the
\emph{rejected} completion, used only by the DPO arm, is a strong-LLM caricature of
the same politician answering the same cue. The context is abridged here to two turns
of the four in the original (579 characters).

\begin{dexample}
\dexlabel{prompt (abridged)}
あなたは日本の国会議員「塩川鉄也（日本共産党）」です。国会審議の以下のやり取りに続けて、あなた自身の発言をしてください。\par
\par
塩川鉄也（日本共産党）：今お話しのように、各都道府県の地域別最低賃金が改定されることに留意し、当該最低賃金を下回らないよう適切に対処してください、こういう通知を出しているわけです。
\par
\par
　これは実際に、じゃ、最賃を下回るような実態があるということも言えるんじゃないですか。\par
\par
［…以下、同じ審議の続きを省略…］\par
\par
塩川鉄也（日本共産党）：

\dexgloss{You are the Japanese Diet member “Shiokawa Tetsuya (Japanese Communist Party).” Continue the following exchange in a Diet deliberation with your own remarks. \quad \textsc{Shiokawa}: As you say, the notice instructs ministries to note that each prefecture's regional minimum wage is being revised and to respond so as not to fall below it. Doesn't that itself indicate that there are cases actually falling below the minimum wage? \quad \textsc{Matsuo (National Personnel Authority)}: The Authority has issued the notice and is urging each ministry to take care in applying it, but we do not have a grasp of any actual instances of non-regular staff being paid below the minimum wage. \quad \textsc{Shiokawa}: So isn't grasping the actual situation and correcting it precisely your job? [\ldots] \quad \textsc{Matsuo}: National public employees are not subject to the Minimum Wage Act, but we have issued a notice and have been guiding each ministry towards appropriate pay, and intend to continue doing so.}
\dexlabel{chosen: the real utterance}
ですから、この間、最低賃金の改定が行われているわけですから、実態が本当にどうなっているのかということをきちんと確認するという、そういうタイミングにもあるんじゃないですか。
\par
　各府省における非常勤の最低賃金がどういう実態なのか、これをしっかりと調べるということが、今全体として最低賃金を引き上げようという取組が行われている中では極めて重要だと思うんですが、そういう実態調査を行う、人事院総裁、いかがですか。

\dexgloss{So, given that revisions to the minimum wage have been taking place, isn't this precisely the moment to confirm properly what the actual situation is? At a time when efforts are being made to raise the minimum wage across the board, I think it is extremely important to investigate thoroughly what the actual state of non-regular minimum wages is in each ministry. Will you conduct such a survey of actual conditions? President of the National Personnel Authority, what do you say?}
\dexlabel{rejected: machine-generated caricature (gemini-2.5-flash-lite), not a real statement}
最低賃金法が適用されないからといって、実質的に最低賃金以下になってしまっている実態があるのであれば、それは問題であるはずです。把握していないというだけで、実態がないとは断定できないのではないでしょうか。是正指導を行うのであれば、まずは実態を正確に把握するところから始めるべきです。

\dexgloss{Even if the Minimum Wage Act does not apply, if there is a reality in which pay effectively falls below the minimum wage, that ought to be a problem. Merely saying you have not grasped it does not permit the conclusion that no such reality exists. If corrective guidance is to be given, one should begin by grasping the actual situation accurately.}
\end{dexample}

\paragraph{The content of the preference contrast.} The caricature is not an
incompetent generation: it is fluent, on topic, and argues the same side of the
question. What it lacks is \emph{procedure}. The real utterance names the official it
addresses, situates the demand within an ongoing policy movement, and closes with a
specific institutional request---that a survey of actual conditions be conducted. The
caricature closes with a generalisation. This is the distinction the preference
objective is constructed to capture, and Section~\ref{sec:objective} establishes that
capturing it does not require the policy to move toward the politician's own
distribution: any direction that separates the real utterance from this caricature
suffices, and most such directions are uninformative about the speaker's identity.
\subsection{A sparse-target instance}

Target \textbf{齋藤裕喜} (立憲民主党・無所属),
sitting of 2024-12-19, speech \texttt{121604006X00220241219\_066} --- a member with
82 speech records, in the second-sparsest band of our sample. Targets carry
prompt/completion pairs only: a held-out target is never given an adapter, so no
\emph{rejected} side is generated for it. Note the difference in register that the
sampling design was intended to expose: this is a committee question closing with an
explicit request for an answer, rather than the prepared set-piece characteristic of a
party leader.

\begin{dexample}
\dexlabel{prompt (abridged)}
あなたは日本の国会議員「齋藤裕喜（立憲民主党・無所属）」です。国会審議の以下のやり取りに続けて、あなた自身の発言をしてください。\par
\par
浅尾慶一郎（自由民主党）：特定帰還居住区域については、帰還の意向のある全ての住民の方々が帰還できるよう、昨年、制度が創設されたところであります。この制度においては、当該区域を含む計画を自治体が申請し、国による認定がされるものであります。
\par
\par
［…以下、同じ審議の続きを省略…］\par
\par
齋藤裕喜（立憲民主党・無所属）：

\dexgloss{You are the Japanese Diet member “Saito Yuki (Constitutional Democratic Party / Independents).” Continue the following exchange with your own remarks. \quad \textsc{Asao (Liberal Democratic Party, Environment Minister)}: [\ldots] The Ministry of the Environment will steadily carry out the decontamination needed so that residents who wish to return may do so at the earliest possible date. Since decontamination aims to reduce effects on human health and the living environment, for forest land we are decontaminating areas adjacent to residential land. [\ldots] Under the policy of eventually lifting evacuation orders for all difficult-to-return zones, the government as a whole will address the matter.}
\dexlabel{chosen: the real utterance}
次に、福島県や近隣県の放射線モニタリングと維持管理についてお尋ねいたします。
\par
　十三年九か月が経過いたしましたが、福島県や近隣県、原子力発災当時の全国各地の状況について、この間、生活圏や農林水産業の風評の被害もまだ払拭されていない状況です。住民の方々や生産者は非常に困っております。
\par
　空間線量だけではなく、住環境の放射線の継続的な維持管理、福島県、近隣県における農地のモニタリング、生活圏の地表面や農林水産物の放射線量を継続的に測ることが、透明な情報を開示する上で、そして風評被害を防ぎ、安心にもつながると思いますが、いかがでしょうか。御答弁よろしくお願いいたします。

\dexgloss{Next, I should like to ask about radiation monitoring and its continued management in Fukushima Prefecture and neighbouring prefectures. Thirteen years and nine months have passed, yet reputational damage to living areas and to agriculture, forestry and fisheries has still not been dispelled, and residents and producers are in great difficulty. Not only ambient dose rates, but continued management of radiation in the living environment, monitoring of farmland, and continuous measurement of ground-surface and agricultural product radiation levels would, I believe, serve transparent disclosure, prevent reputational damage and provide reassurance. What is your view? Your answer, please.}
\end{dexample}
\subsection{The in-context prefix (baseline 4)}

\texttt{ICL} prepends the target's own dev-fold utterances --- the \emph{chosen} sides,
not their debate contexts, which would exhaust the context window without saying
anything more about how the politician speaks --- to every held-out prompt. The block
is capped at 12k characters, which for a 24-instance budget is a prefix of roughly 9k
tokens paid at every inference. Two demonstrations are shown, each clipped after 150
characters; the run uses every dev instance in full.

\begin{dexample}
\dexlabel{prefix template, with two of the demonstrations}
以下は、あなたがこれまでの国会審議で実際に行った発言です。\par
\par
・立憲民主党の齋藤裕喜と申します。
\par
　本日、厚生労働委員会では初めての質問をさせていただきたいと思います。上野大臣、政府の方々、どうぞよろしくお願いいたします。
\par
　まず初めに、お手元にまだ資料は届いていないですね。お配りしている資料についてちょっと、国の財政の視点、医療や薬局、そして家計の視点、こ［…］\par
\par
・ありがとうございます。
\par
　確かに、財源の問題で全世代型社会保障ということをおっしゃっているのは承知しております。
\par
　誰しもがやはり老いて、人生の最期を迎えていくわけですね。そんな中、病気と闘っている方々、難病を患っている方々、治療を継続して高額になっている方々の支援の拡充をこの財政の中でもしてい［…］\par
\par
［…残りの発言を省略…］\par
\par
上記の発言と同じ立場・話し方で、次の審議に応じてください。

\dexgloss{Below are remarks you have actually made in past Diet deliberations. [\ldots] Respond to the following deliberation in the same position and manner of speaking as the remarks above.}
\end{dexample}
\subsection{The UTAS elicitation prompt (Section~\ref{sec:utas})}

Numeric-label scoring: the options are enumerated in the prompt and only the single
label token the model would emit is scored, so every candidate continuation is exactly
one token and the answer distribution is length-unbiased. Item
\texttt{Q4\_1} of the 2024 House of Representatives wave is shown. The
persona is conditioned on the name alone --- no speech, no adapter and no merge is
required to run this, which is what made the 594-persona population control in
Section~\ref{sec:utas} possible.

\begin{dexample}
\dexlabel{system}
あなたは日本の国会議員、高市早苗です。以下の政策に関する意見について、選択肢の中からあなたの立場に最も近いものを一つだけ答えてください。

\dexgloss{You are a member of the Japanese Diet, Takaichi Sanae. For the following policy opinion, answer with the single option closest to your own position.}
\dexlabel{user}
日本の防衛能力は強化されるべきである。\par
1. そう思う\par
2. どちらかといえばそう思う\par
3. どちらともいえない\par
4. どちらかといえばそう思わない\par
5. そう思わない\par
\par
上記の選択肢の中から、あなたの立場に最も近いものの番号（1〜5）だけを答えてください。

\dexgloss{“Japan's defense capabilities should be strengthened.” \quad 1. Agree \quad 2. Somewhat agree \quad 3. Neither \quad 4. Somewhat disagree \quad 5. Disagree \quad Answer only with the number of the option closest to your position.}
\dexlabel{scored continuation}
1 / 2 / 3 / 4 / 5

\dexgloss{One token. The reported expectation is $\sum_k k\,p(k)$ on the 1--5 scale; the rank metric uses the resulting 33-vector.}
\end{dexample}

The item wording here is a machine translation into Japanese of the English codebook,
performed once and reused for every persona and configuration, so no configuration is
advantaged by it.
\subsection{Instance counts and lengths}

Character counts rather than token counts, because the export budgets context in
characters (Japanese runs at roughly one token per character for this backbone). The
prompt cap is what keeps prompt-plus-completion inside the trainer's 1\,024-token
limit.

\begin{table}[t]
\centering
\caption{Anchor training sets. Each anchor contributes 1\,500 instances; the
\emph{rejected} column exists only in the DPO arm. Medians.}
\label{tab:anchorstats}
\begin{tabular}{llrrrrr}
\toprule
Anchor & Party & instances & prompt & turns & chosen & rejected \\
\midrule
岸田文雄 & LDP & 1,500 & 1008 & 4 & 303 & 208 \\
塩川鉄也 & JCP & 1,500 & 937 & 5 & 275 & 191 \\
福島みずほ & SDP & 1,500 & 920 & 5 & 196 & 184 \\
上田清司 & DPJ & 1,500 & 899 & 5 & 285 & 178 \\
\bottomrule
\end{tabular}
\end{table}

\begin{table}[t]
\centering
\caption{Held-out targets by prominence band: median instances available per target
and median lengths in characters. Only 30 instances are ever used, whatever a target
has; the export caps extraction at 300.}
\label{tab:targetstats}
\begin{tabular}{lrrrrr}
\toprule
Band (speech records) & targets & instances & prompt & turns & chosen \\
\midrule
originals & 3 & 300 & 926 & 4 & 318 \\
3000+ & 5 & 300 & 960 & 4 & 266 \\
1000--3000 & 5 & 300 & 919 & 4 & 286 \\
500--1000 & 5 & 300 & 989 & 4 & 365 \\
250--500 & 5 & 300 & 966 & 3 & 408 \\
120--250 & 5 & 162 & 1013 & 4 & 398 \\
60--120 & 5 & 102 & 985 & 3 & 448 \\
\bottomrule
\end{tabular}
\end{table}

Two things in Table~\ref{tab:targetstats} are worth noting. The 300-instance extraction
cap binds for every band except the two sparsest, so for most targets the binding
constraint on the experiment is the 30-utterance \emph{budget} rather than data
availability --- which is the sense in which the method is low-resource by construction
rather than by circumstance. Second, and less expected, \textbf{sparse targets speak in
longer turns}: median completion length rises from 266 characters in the 3000+ band to 448 in the 60--120 band, a 1.7$\times$ difference in the opposite direction to speech count. Frequent speakers accumulate a
record through many short interventions, whereas a backbencher's record is a handful of
long committee questions. This is a plausible part of why the margins in
Section~\ref{sec:prominence} do not decay with prominence --- a fixed 30-utterance
budget buys more text on a sparse target than the speech count suggests --- and it is a
caveat against reading speech-record counts as a proxy for available training signal.
